\documentclass[10pt,a4paper]{amsart}
\usepackage[margin=19mm]{geometry}
\usepackage{amsmath,amssymb,mathtools}
\usepackage[scr=rsfs,cal=euler]{mathalfa}
\usepackage{xfrac}
\usepackage[unicode,hidelinks,bookmarks=false]{hyperref}
\newcommand{\ie}{i.e.\ }
\newcommand{\cf}{cf.\ }
\newcommand{\resp}{respectively\ }
\newcommand{\Subsection}{\subsection}
\providecommand{\nobreakdash}{\nobreak\hbox{-}}
\setlength{\emergencystretch}{2em}
\multlinegap0pt
\usepackage{amssymb}
\usepackage{mathtools}
\usepackage{bm}
\usepackage[shortlabels]{enumitem}
\newtheorem{lemma}{Lemma}
\newtheorem{theorem}{Theorem}
\usepackage{cleveref}
\crefname{lemma}{Lemma}{Lemmas}
\crefname{theorem}{Theorem}{Theorems}
\newcommand{\Qhat}{\widehat Q}\newcommand{\Qp}{\widehat{\mathfrak Q}}
\newcommand{\R}{\mathbb R}\newcommand{\dd}{\,d}\newcommand{\HH}{\mathcal H}
\title{\textbf{Threshold Decomposition of Vorticity Stretching and Palinstrophy}}
\author{Rishad Shahmurov}
\date{Source: arXiv:2605.01873v3 (CC BY 4.0)}
\begin{document}
\maketitle

\noindent\emph{Source and licence.} \textbf{Threshold Decomposition of Vorticity Stretching and Palinstrophy}. Version arXiv:2605.01873v3 (CC BY 4.0), distributed by arXiv under the Creative Commons Attribution 4.0 International licence, http://creativecommons.org/licenses/by/4.0/, as recorded in the arXiv licence metadata for this exact version and shown on the abstract page https://arxiv.org/abs/2605.01873v3. Full licence notice: \texttt{licenses/arxiv-2605.01873-CC-BY-4.0.txt}.\par\smallskip

\noindent\textit{Editorial scope.} Counted unit: the article's theorem thm:threshold (source0.tex lines 177--191). The article's complete original proof of it is delivered with it (source0.tex lines 192--325, one \texttt{proof} environment of 134 lines). No mathematics is written, repaired or re-derived: the statement and the proof are the article's own lines, in the article's own notation, and no retained line is substituted or re-flowed.  Retained uncounted as the source-local context the proof needs: lines 88--94.  Nothing was omitted from any retained span, so the package carries no omission record.  No retained line contains a citation, so the wrapper carries no bibliography and no bibliography entry.  No retained line contains a figure environment, an image-inclusion command or drawing code, so no image is delivered and the package declares no asset dependency.  Editorial formatting only, in addition: an AMS \texttt{amsart} preamble that restates the article's own declarations and macros used by the retained lines, namely the environments \texttt{lemma}, \texttt{theorem} on separate counters, together with the standard packages those lines need.  The retained environments print in this document as Lemma 1, Theorem 1 in order of appearance, whereas the article prints the same items with the numbering of its own full document.  The complete native source of the article as distributed in the arXiv e-print for this exact version is bundled unchanged as \texttt{sources/199667/source0.tex}.
\begin{lemma}\label{lem:magnitude}
On $\{f>0\}$,
\[
(\partial_t+u\cdot\nabla-\nu\Delta)f
=\alpha f-\nu f|\nabla\xi|^2.
\]
\end{lemma}

\begin{theorem}\label{thm:threshold}
For almost every $a>0$, define
\[
\Qhat_I(a)=\int_I\int_{\{f>a\}}\alpha f\,\dd x\,\dd t
-\left[\int_{\R^3}(f-a)_+\,\dd x\right]_{t_0}^{t_1}.
\]
Then
\[
\boxed{
\Qhat_I(a)=
\nu\int_I\int_{\{f>a\}}f|\nabla\xi|^2\,\dd x\,\dd t
+\nu\int_I\int_{\{f=a\}}|\nabla f|\,\dd\HH^2\,\dd t\ge0.
}
\]
\end{theorem}

\begin{proof}
Fix $a>0$.  We first derive a smooth identity and only then pass to the
nonsmooth truncation.  This avoids any ambiguity at points where $\omega=0$.

Choose a nonnegative even mollifier $\rho\in C_c^\infty((-1,1))$ with
$\int_\R\rho=1$, and set
\[
\rho_\varepsilon(s)=\varepsilon^{-1}\rho(s/\varepsilon).
\]
Let $H_\varepsilon$ be the smooth increasing function satisfying
\[
H_\varepsilon'=\rho_\varepsilon,
\qquad
0\le H_\varepsilon\le1,
\qquad
H_\varepsilon(s)=0\ \text{for }s\le-\varepsilon,
\qquad
H_\varepsilon(s)=1\ \text{for }s\ge\varepsilon.
\]
Define
\[
\Phi_{a,\varepsilon}(s)=\int_0^sH_\varepsilon(\sigma-a)\,\dd\sigma.
\]
Then $\Phi_{a,\varepsilon}$ is convex,
\[
\Phi_{a,\varepsilon}'(s)=H_\varepsilon(s-a),
\qquad
\Phi_{a,\varepsilon}''(s)=\rho_\varepsilon(s-a)\ge0,
\]
and $\Phi_{a,\varepsilon}(s)\to(s-a)_+$ locally uniformly as
$\varepsilon\downarrow0$.  We may take $0<\varepsilon<a/2$.  On the support
of $\Phi'_{a,\varepsilon}(f)$ one then has $f>a/2>0$, so
\cref{lem:magnitude} is valid pointwise wherever the multiplier is nonzero.

Multiply that identity by $\Phi'_{a,\varepsilon}(f)$ and integrate over
$\R^3$.  The chain rule gives
\[
\Phi'_{a,\varepsilon}(f)(\partial_t+u\cdot\nabla)f
=(\partial_t+u\cdot\nabla)\Phi_{a,\varepsilon}(f).
\]
Because $\nabla\cdot u=0$ and the solution has the stated decay,
\[
\int_{\R^3}u\cdot\nabla\Phi_{a,\varepsilon}(f)\,\dd x
=\int_{\R^3}\nabla\cdot\bigl(u\Phi_{a,\varepsilon}(f)\bigr)\,\dd x=0.
\]
For the Laplacian term, integration by parts and the chain rule give
\[
\begin{aligned}
\int_{\R^3}\Phi'_{a,\varepsilon}(f)\Delta f\,\dd x
&=-\int_{\R^3}\nabla\bigl(\Phi'_{a,\varepsilon}(f)\bigr)\cdot\nabla f\,\dd x\\
&=-\int_{\R^3}\Phi''_{a,\varepsilon}(f)|\nabla f|^2\,\dd x.
\end{aligned}
\]
We therefore obtain, for every such $\varepsilon$,
\begin{equation}\label{eq:smooth-threshold}
\begin{aligned}
\frac{\dd}{\dd t}\int_{\R^3}\Phi_{a,\varepsilon}(f)\,\dd x
={}&\int_{\R^3}\Phi'_{a,\varepsilon}(f)\alpha f\,\dd x\\
&-\nu\int_{\R^3}\Phi'_{a,\varepsilon}(f)f|\nabla\xi|^2\,\dd x\\
&-\nu\int_{\R^3}\rho_\varepsilon(f-a)|\nabla f|^2\,\dd x.
\end{aligned}
\end{equation}
Integrating from $t_0$ to $t_1$ produces the finite-interval identity with
smooth truncation.

We now pass to the limit.  First,
\[
\Phi_{a,\varepsilon}(f)\longrightarrow(f-a)_+,
\qquad
\Phi'_{a,\varepsilon}(f)\longrightarrow\mathbf1_{\{f>a\}}
\]
at every point with $f\ne a$.  For almost every $a$, the spacetime level set
$\{(x,t):f(x,t)=a\}$ has four-dimensional Lebesgue measure zero.  Indeed,
this follows from Fubini together with the elementary fact that a real-valued
measurable function has positive-measure level sets for at most countably
many levels on each finite-measure truncation of spacetime.  The smoothness
and decay assumptions give an integrable majorant for the stretching and
bulk-direction terms on the finite interval, so dominated convergence yields
\[
\int_I\!\int\Phi'_{a,\varepsilon}(f)\alpha f
\longrightarrow
\int_I\!\int_{\{f>a\}}\alpha f
\]
and
\[
\int_I\!\int\Phi'_{a,\varepsilon}(f)f|\nabla\xi|^2
\longrightarrow
\int_I\!\int_{\{f>a\}}f|\nabla\xi|^2.
\]
The endpoint term converges to
\[
\left[\int_{\R^3}(f-a)_+\,\dd x\right]_{t_0}^{t_1}.
\]

It remains to identify the last term in \eqref{eq:smooth-threshold}.  Define
for almost every $s>0$
\[
G(s)=\int_I\int_{\{f(\cdot,t)=s\}}|\nabla f|\,\dd\HH^2\,\dd t.
\]
The coarea formula, first in space and then integrated in time, gives for
any compact interval $[c,d]\subset(0,\infty)$
\[
\int_c^dG(s)\,\dd s
=\int_I\int_{\{c<f<d\}}|\nabla f|^2\,\dd x\,\dd t<\infty.
\]
Hence $G\in L^1_{\rm loc}(0,\infty)$.  A second application of coarea gives
\[
\begin{aligned}
\int_I\int_{\R^3}\rho_\varepsilon(f-a)|\nabla f|^2\,\dd x\,\dd t
&=\int_0^\infty\rho_\varepsilon(s-a)G(s)\,\dd s.
\end{aligned}
\]
The right-hand side is an approximate identity applied to $G$.  By the
Lebesgue differentiation theorem it converges to $G(a)$ for almost every
$a>0$.  Thus
\[
\lim_{\varepsilon\downarrow0}
\int_I\int\rho_\varepsilon(f-a)|\nabla f|^2
=
\int_I\int_{\{f=a\}}|\nabla f|\,\dd\HH^2\,\dd t
\]
for almost every $a$.

Passing to the limit in the time-integrated form of
\eqref{eq:smooth-threshold} and moving the endpoint increment to the left gives
exactly
\[
\Qhat_I(a)=
\nu\int_I\int_{\{f>a\}}f|\nabla\xi|^2
+\nu\int_I\int_{\{f=a\}}|\nabla f|\,\dd\HH^2\,\dd t.
\]
Both terms on the right are nonnegative, proving the final inequality as
well.
\end{proof}

\end{document}
